\chapter{Marco tecnológico: TileDB en detalle}\label{anexo-h.-marco-tecnologico-tiledb}

Este anexo amplía el marco tecnológico presentado en el capítulo~\ref{analisis-de-tecnologias}, centrándose en los aspectos internos de TileDB que explican las decisiones de diseño del sistema. No pretende sustituir a la documentación oficial \cite{tiledb_web}, sino reunir en un solo lugar los conceptos necesarios para entender por qué el almacén se comporta como lo hace.

\section{Modelo de datos: arrays densos y dispersos}\label{h-modelo}

TileDB organiza los datos como arrays multidimensionales. Un array se define por dos elementos: las \textbf{dimensiones}, que forman el sistema de coordenadas por el que se indexa (en este trabajo, tiempo, fila y columna), y los \textbf{atributos}, que son los valores almacenados en cada celda (las variables hidráulicas H, QX, QY y MH). La diferencia con una imagen ráster convencional es que las dimensiones no son solo ejes de una matriz: son un índice por el que se puede consultar de forma eficiente.

Un array puede ser \textbf{denso} (se almacenan todas las celdas, haya o no dato) o \textbf{disperso} (\textit{sparse}, solo se almacenan las celdas con dato, guardando explícitamente sus coordenadas). La elección es la decisión central de este trabajo: en una inundación, en cada instante solo alrededor del 8\,\% del dominio está mojado, de modo que un array disperso almacena únicamente esas celdas y evita el coste de las celdas secas, mientras que uno denso las guardaría todas.

Conviene matizar que el modelo disperso no es superior en términos absolutos. Guardar explícitamente las coordenadas de cada celda añade un sobrecoste que solo compensa cuando los datos son realmente dispersos; sobre un dominio mayoritariamente lleno, un array denso ocupa menos y se recorre de forma más directa. Existe además un compromiso (\textit{trade-off}) entre coste de almacenamiento y coste de procesamiento: según cómo esté implementado el procesador de consultas, recuperar y recombinar las coordenadas explícitas de un array disperso puede tener un coste de cálculo mayor que recorrer una matriz densa contigua. En este trabajo la dispersión es muy alta (en torno al 92\,\% de celdas secas) y los patrones de acceso son por ventana, dos condiciones que inclinan claramente el balance hacia el array disperso; pero la decisión correcta depende siempre de las características concretas de los datos y de la carga de consultas.

\section{Modelo de escritura: fragmentos y consolidación}\label{h-fragmentos}

TileDB no modifica los datos en el sitio. Cada operación de escritura crea un \textbf{fragmento} (\textit{fragment}) independiente e inmutable, una capa con los datos de esa escritura, y una lectura posterior combina de forma transparente todos los fragmentos para devolver la vista completa del array. Este modelo, inspirado en las estructuras de tipo \textit{log-structured} habituales en los sistemas de almacenamiento modernos \cite{oneil_lsm}, favorece la escritura concurrente y evita bloqueos.

El proceso ETL de este trabajo escribe los datos paso a paso, lo que genera \textbf{20 fragmentos}, uno por instante temporal. Esta estructura no es casual: como cada fragmento contiene un único paso, las consultas que necesitan un solo instante leen exactamente un fragmento, y el comportamiento del sistema se mantiene predecible.

Frente a este reparto, TileDB ofrece una operación de \textbf{consolidación} que fusiona varios fragmentos en uno solo, pensada para escenarios con miles de escrituras pequeñas. En este sistema, sin embargo, consolidar los 20 fragmentos en uno resulta contraproducente: al mezclar los pasos temporales en tiles compartidos, una consulta que solo necesita un instante se ve obligada a leer y descomprimir tiles que contienen también otros pasos. Este fenómeno, conocido como \textbf{amplificación de tiles} (\textit{tile amplification}, leer más datos de los necesarios), llegó a multiplicar por 16 el tiempo de las consultas multi-paso durante el desarrollo. Por ello el sistema mantiene deliberadamente los 20 fragmentos sin consolidar: que existan 20 fragmentos es el estado correcto del array, no un defecto.

\section{Organización física: tiles, índice y lectura por ventana}\label{h-tiles}

Dentro de un array, la unidad física de almacenamiento, compresión y lectura es el \textbf{tile}. Un tile agrupa un bloque de celdas que se comprime y se lee como una sola pieza; en este sistema los tiles espaciales son de 256$\times$256 celdas \cite{tiledb_datalayout}. La consecuencia práctica es que la granularidad de acceso es el tile: para leer una sola celda hay que descomprimir el tile completo que la contiene.

La compresión de cada tile se configura como una \textbf{cadena de filtros} (\textit{filter pipeline}) que se aplica antes de escribirlo \cite{tiledb_compression}. El sistema emplea dos filtros encadenados: primero \textit{ByteShuffle}, que reordena los bytes de los valores \texttt{float32} agrupando los de igual posición (todos los primeros bytes juntos, luego todos los segundos, y así sucesivamente), y después ZSTD, que comprime el resultado. El reordenamiento previo es clave porque los bytes de igual posición de valores hidráulicos cercanos son muy parecidos, lo que crea largas secuencias repetidas que ZSTD comprime mucho mejor. La elección concreta de filtro y nivel se determinó experimentalmente en el \textit{benchmark} de almacenamiento del Anexo~\ref{anexo-b.-tablas-de-resultados-completas}.

Para localizar los tiles relevantes sin recorrer todo el array, TileDB ordena las celdas de un array disperso según un \textbf{orden global} sobre sus coordenadas y mantiene un índice (de tipo R-tree) que indica con rapidez qué tiles intersecan una región consultada \cite{tiledb_indexing}. La combinación del modelo disperso, este índice y los tiles habilita la \textbf{lectura por ventana} (\textit{windowed read}), ilustrada en la Figura~\ref{fig:windowed-read}: dada una región del dominio, TileDB consulta el índice para identificar los tiles que la intersecan, lee solo esos tiles y descomprime únicamente lo necesario. Esta es la base de la ventaja de rendimiento medida frente a GeoTIFF (capítulo~\ref{validacion-del-sistema}): mientras que leer una subregión de un GeoTIFF clásico obliga a recorrer estructuras pensadas para el acceso secuencial, el índice disperso de TileDB acota el trabajo a la zona consultada.

\begin{figure}[ht]
\centering
\begin{tikzpicture}[font=\small\sffamily, x=0.8cm, y=0.8cm]
  % --- Dominio: fondo seco ---
  \fill[gray!7] (0,0) rectangle (12,9);

  % --- Tiles que interseca la ventana (se leen enteros) ---
  \foreach \tx/\ty in {6/3, 9/3, 6/6, 9/6} {
    \fill[orange!25] (\tx,\ty) rectangle ++(3,3);
  }

  % --- Celdas humedas que NO se leen (su tile no toca la ventana) ---
  \foreach \x/\y in {0/1,1/2,2/1,0/5,1/4,3/3,4/5,1/7,2/6,3/8,4/8,0/8,5/1,5/6,7/1,9/0,10/2,11/1,6/1,8/0} {
    \fill[blue!22] (\x,\y) rectangle ++(1,1);
  }

  % --- Celdas humedas que SI se leen (dentro de los tiles naranjas) ---
  \foreach \x/\y in {7/4,8/5,6/3,10/4,9/5,7/7,8/6,10/8,11/6,9/7,6/5,11/4} {
    \fill[blue!72] (\x,\y) rectangle ++(1,1);
  }

  % --- Rejilla fina de celdas ---
  \draw[gray!18, very thin] (0,0) grid (12,9);

  % --- Rejilla gruesa de tiles (3x3 celdas) ---
  \foreach \i in {0,3,...,12} { \draw[gray!60, thick] (\i,0)--(\i,9); }
  \foreach \j in {0,3,...,9}  { \draw[gray!60, thick] (0,\j)--(12,\j); }

  % --- Ventana de consulta (no alineada con los tiles, a proposito) ---
  \draw[red!75!black, very thick, dashed] (6.4,3.4) rectangle (11.4,8.4);
  \node[red!75!black, font=\scriptsize\bfseries, anchor=south west]
        at (6.45,8.55) {ventana de consulta};

  % --- Etiqueta: un tile ---
  \draw[gray!75, -{Latex[length=1.6mm]}] (1.5,-0.55) -- (2.4,1.4);
  \node[gray!55!black, font=\scriptsize, anchor=north] at (1.5,-0.55)
        {1 tile = 256$\times$256 celdas};

  % ===================== LEYENDA =====================
  \begin{scope}[shift={(0,-2.2)}]
    % columna 1
    \fill[gray!7]    (0,-0.0) rectangle ++(0.7,0.7); \draw[gray!60] (0,-0.0) rectangle ++(0.7,0.7);
    \node[anchor=west, font=\scriptsize] at (0.85,0.35) {dominio seco (no se lee)};
    \fill[orange!25] (0,-1.0) rectangle ++(0.7,0.7); \draw[gray!60] (0,-1.0) rectangle ++(0.7,0.7);
    \node[anchor=west, font=\scriptsize] at (0.85,-0.65) {tile leído y descomprimido};
    \draw[red!75!black, very thick, dashed] (0,-2.0) rectangle ++(0.7,0.7);
    \node[anchor=west, font=\scriptsize] at (0.85,-1.65) {ventana de consulta};
    % columna 2
    \fill[blue!72] (6.6,-0.0) rectangle ++(0.7,0.7); \draw[gray!50] (6.6,-0.0) rectangle ++(0.7,0.7);
    \node[anchor=west, font=\scriptsize] at (7.45,0.35) {celda húmeda leída};
    \fill[blue!22] (6.6,-1.0) rectangle ++(0.7,0.7); \draw[gray!50] (6.6,-1.0) rectangle ++(0.7,0.7);
    \node[anchor=west, font=\scriptsize] at (7.45,-0.65) {celda húmeda no leída};
  \end{scope}
\end{tikzpicture}
\caption{Lectura por ventana (\textit{windowed read}) sobre un array disperso. Ante una consulta acotada a una región (recuadro rojo discontinuo), el índice de TileDB identifica los \textbf{tiles} que la intersecan (en naranja) y descomprime únicamente esos, en lugar de recorrer todo el dominio. Solo se leen las celdas húmedas contenidas en esos tiles (azul intenso); el resto de celdas húmedas (azul claro) y la inmensa mayoría de celdas secas permanecen en disco sin descomprimirse. La ventana no coincide con los límites de los tiles a propósito: como el tile es la unidad mínima de lectura, basta con que la ventana toque parte de un tile para leerlo entero.}
\label{fig:windowed-read}
\end{figure}

\section{Encaje con las alternativas}\label{h-alternativas}

Las alternativas valoradas se analizan en el capítulo~\ref{analisis-de-tecnologias}. En resumen, SciDB, la referencia académica de las bases de datos de arrays, carece de versiones públicas recientes de su edición abierta y exige servidor dedicado; las bibliotecas de array sobre fichero (HDF5, Zarr) no ofrecen de serie la indexación dispersa sobre coordenadas que aquí resulta decisiva. TileDB combina el modelo disperso, la compresión por filtros encadenados y el acceso por ventana en una biblioteca embebida sin servidor \cite{tiledb_github}, lo que se ajusta a los requisitos del problema sin la complejidad operativa de un sistema cliente-servidor.
